\newcommand{\DoNotLoadEpstopdf}{}
\documentclass[11pt,reqno]{amsart}
\usepackage[T1]{fontenc}
\usepackage{lmodern,amsmath,amssymb,amsthm,mathtools,mathrsfs,microtype}
\usepackage[margin=1.1in]{geometry}
\usepackage{xr-hyper}
\usepackage[hidelinks]{hyperref}
\externaldocument[supp-]{build/supplement}[supplement.pdf]
\hypersetup{pdftitle={Collision records, arithmetic, and raw local inversion in a triangular Lorentz gas},pdfauthor={Qian Qi}}
\setlength{\emergencystretch}{2em}
\allowdisplaybreaks[2]
\numberwithin{equation}{section}
\newtheorem{theorem}{Theorem}[section]
\newtheorem{lemma}[theorem]{Lemma}
\newtheorem{proposition}[theorem]{Proposition}
\newtheorem{corollary}[theorem]{Corollary}
\theoremstyle{definition}
\newtheorem{definition}[theorem]{Definition}
\newtheorem{inputstatement}[theorem]{Input statement}
\theoremstyle{remark}\newtheorem{remark}[theorem]{Remark}
\newcommand{\R}{\mathbb R}\newcommand{\Z}{\mathbb Z}
\newcommand{\one}{\mathbf1}\newcommand{\dd}{\,\mathrm d}
\newcommand{\esssup}{\operatorname*{ess\,sup}}
\newcommand{\diver}{\operatorname{div}}
\newcommand{\TV}{\operatorname{TV}}
\title[Collision records and raw local inversion]{Collision records, arithmetic, and raw local inversion in a triangular Lorentz gas}
\author{Qian Qi}
\date{October 11, 2026. A2-DYN, revision 83}
\subjclass[2020]{37D50, 37A50, 60F05, 28A75}
\keywords{Sinai billiards, exact collision records, local limits, fractional optical potentials, first-exit surface sources, exact-label coarea transport}
\begin{document}
\raggedbottom
\begin{abstract}
We study the pointwise local limit problem for the full displacement,
collision-count, occupation and roof record of actual returns in a
triangular Lorentz gas. Actual fractional inverse-gradient insertions
yield complete optical transport estimates without an initial critical
cutoff. The optical loss and genuine grazing starting mass obey local
weak-$L^q$ bounds for every fixed $2<q<4$, with an additional power
on the critical initial source. We estimate the actual nongrazing
first-exit surface assignment by paying both signs of optical variation.
Together with all-depth physical margin localization, this gives a
collision-normalized $O(r^{33/32})$ bound for its mass on each previously
fixed roof interval of length $r$. The exact labels, original guard and
one-sided first-exit selection are retained. A separate ordered-tail
argument identifies the microscopic estimate still required for the
original essential-height law; local norm and surface mass bounds are
not substituted for that pointwise estimate.
\end{abstract}
\maketitle

\section{The unchanged physical problem}
Put $R_-=9/20$ and $R_+=47/100$, and let $R\in I=[R_-,R_+]$. Circular scatterers have centers in
$\Lambda=\Z(1,0)+\Z(1/2,\sqrt3/2)$. Write $T_R$ for the regular
collision map and $(\alpha,p)\in(\R/2\pi\Z)\times(-1,1)$ for its
coordinates; $p$ is the outgoing tangential velocity and
$\dd\nu=\dd\alpha\,\dd p/(4\pi)$. Physical flights satisfy
$\tau_-\le\tau_R\le\tau_+$, where $\tau_-=3/50$ and $\tau_+$
is the fixed finite-horizon bound.

The section is specified here to fix every event in the argument.
Put $h_0=1/200$. Let $Y_R$ be the union of the four open coordinate
squares of half-side $h_0$ centered at
\[
 (0,0),\quad(\pi/6,-1/2),\quad
 (-\theta_R,\sin\theta_R),\quad
 (\pi/3-\theta_R,\sin\theta_R).
\]
Here $\theta_R=2\arctan q_R$, and $q_R$ is the unique root in
$(49/100,53/100)$ of
\[
 8Rq(1+q^2)-4q(1-q^2)+\sqrt3(1-6q^2+q^4)=0.
\]
Set
\[
 U=\{|\alpha-\pi/6|<3/50,\ |p|<3/20\},\qquad
 Y_R^*=Y_R\cup U,\qquad c_*=\frac{91}{10000\pi}.
\]
The rectangles are disjoint and $\nu(Y_R^*)=c_*$. We always use
$\nu_R^*=c_*^{-1}\nu|_{Y_R^*}$, not a renormalization of a smaller
source.

At $m$ flights, let $L=L_{m,R}$ be their total length,
$K^{\rm c}_{m,R}$ their lattice displacement, and
$A_{m,R}=\sum_{j=0}^{m-1}\one_{Y_R^*}\circ T_R^j$. The event is
\begin{equation}\label{eq:v76-event}
 \Pi_{n,k,m,R}=Y_R^*\cap\{K^{\rm c}_{m,R}=k,\ A_{m,R}=n\}
                         \cap T_R^{-m}Y_R^*.
\end{equation}
The initial visit is counted; the terminal visit is tested separately.
On this event the actual $n$th return has collision count $m$ and
roof $L$. Denote its density by $p_{n,R}(k,m,t)$.

Write $c_j=\sqrt{1-p_j^2}$ and let $\Delta_j$ be the least clearance
of flight $j$ from any third disk, excluding its two incident disks.
Keep the original smooth function $0\le\Theta\le1$, zero on
$(-\infty,1]$ and one on $[2,\infty)$. With $q_0=47/53$,
$d_j=\min(j,m-j)$ and $d'_j=\min(j,m-1-j)$, set
\begin{align}
 H^\varepsilon&=
 \prod_{j=0}^{m}\Theta\!\left(\frac{c_j}{\varepsilon q_0^{d_j/4}}\right)
 \prod_{j=0}^{m-1}\Theta\!\left(\frac{\Delta_j}{\varepsilon q_0^{d'_j/4}}\right),
 \qquad W^\varepsilon=1-H^\varepsilon,\label{eq:v76-guard}\\
 b^\varepsilon(t)\dd t
   &=L_\#(W^\varepsilon\one_{\Pi_{n,k,m,R}}\nu_R^*).
 \label{eq:v76-source}
\end{align}
In particular first incidence and first clearance outside any selected
critical neighborhoods are still present. No section-distance factor
has been inserted into this guard.

\subsection{The arithmetic reference and the target}
For $y=(k_1,k_2,t,n)$ put
\[
 \mu_R=(0,0,\bar\tau_R,c_*),\quad
 Z=(y-m\mu_R)/\sqrt m,\quad
 \mathcal C_{m,R,n,k}(M)=\{t:|Z|\le M\}.
\]
We retain the finite spectral-peak reference, including peaks close to
but not on the unit circle. To specify its meaning, the collision
record observable is $(\kappa_1,\kappa_2,\tau,\one_{Y_R^*})$.
Let $\lambda_j(R,\vartheta)$ and $\mathsf P_j(R,\vartheta)$ be its
continued simple twisted eigenvalue and projection in each of the
fixed peak charts, and let $\chi_j$ be their frequency partition.
At the real maximum $a_j$ of $\log|\lambda_j|$ put
\[
 z_j=\lambda_j(R,a_j),\quad
 \mu_j=i^{-1}\nabla\log\lambda_j(R,a_j),\quad
 H_j=-\nabla^2\log\lambda_j(R,a_j),
\]
with $\operatorname{Re}H_j\ge aI_4$ for a common $a>0$. If
$\eta_R=\one_{Y_R^*}$ and $\ell$ denotes mass evaluation, put
$B_j=\chi_j(R,a_j)\ell(M_{\eta_R}\mathsf P_j(R,a_j)(\eta_R\nu))$.
Define
\begin{align}
 G_H(z)&=(2\pi)^{-4}\int_{\R^4}e^{-iv\cdot z-v^{\mathsf T}Hv/2}\dd v,
 \notag\\
 \mathcal L_{m,R}(y)&=\operatorname{Re}\left[
 c_*^{-1}\sum_j e^{-ia_j\cdot y}z_j^m B_j
      G_{H_j}\left((y-m\mu_j)/\sqrt m\right)\right].
 \label{eq:v76-reference}
\end{align}
Only the charts active at $R$ are summed. This is the same transition
reference as in the retained spectral theorem. The existence and
uniform estimates for these scalar data belong to that analytic
input, not to the transport argument below. No frequency branch or
arithmetic zero class is deleted.

The governing assertion remains
\begin{equation}\label{eq:v76-target}
 \sup_{R,\,1\le n\le m,\,k\in\Z^2}
 \esssup_{t\in\mathcal C_{m,R,n,k}(M)}
 |m^2p_{n,R}(k,m,t)-\mathcal L_{m,R}(k_1,k_2,t,n)|\longrightarrow0.
\end{equation}
Let $\mathcal E_M$ denote this error's collision limsup. An empty
central target contributes zero. The order of the uniformity is
$\sup_R\esssup_t$, so exceptional density sets may depend on $R$.

\subsection{The complete bounds proved here}
The two positive transport mechanisms will be treated at the level of
the full original source. Let $O_h^\varepsilon$ denote positive optical
loss on survivors, $E_h^{\varepsilon,\mathrm{gr}}$ the starting mass of
genuine grazing first exits, and $\mathcal F_h$ the nongrazing outgoing
first-exit flux assigned to the starting roof. Their precise definitions
are given below and in \eqref{eq:v83-untruncated-flux}. A bar denotes
$h$-averaging over $(0,r)$. For $q>1$ use the normalization
$\mathcal N_q$ in \eqref{eq:v83-lorentz-norm}.

\begin{theorem}[Complete optical and first-exit estimates]
\label{thm:v83-main}
Under the explicit inherited inputs below, for every fixed $2<q<4$
and every previously fixed $r,L>0$,
\begin{align}
 \limsup_m\sup_{R,n,k}\sup_{|J|=L}
 \mathcal N_q\bigl(m^2(\bar O_r^\varepsilon+
             \bar E_r^{\varepsilon,\mathrm{gr}});J\bigr)
                                     &\le C_q(L+r),\label{eq:v83-main-optical}\\
 \limsup_m\sup_{R,n,k}\sup_{|J|=r}
           m^2\int_J\overline{\mathcal F}_r(t)\,dt
                                     &\le C_q r^{33/32}
                               \quad(0<r<r_0).\label{eq:v83-main-surface}
\end{align}
Both sources are untruncated in the initial endpoint gradient and in
all positive internal incidences. They retain the original exact label,
initial and terminal tests, and true one-sided first-exit selection.
\end{theorem}
The first assertion is proved in Theorem~\ref{thm:v83-complete-weak-height}
with $\eta=1-2/q$, and includes sharper bounds on a critical initial
neighborhood. The second is Theorem~\ref{thm:v83-surface-local-mass}.
Their proofs are complete relative to the stated inputs. They do not
assert \eqref{eq:v76-target}, whose remaining essential-height estimate
is formulated in Theorem~\ref{thm:v83-endpoint}. In particular the
window $r$ is never chosen as a function of the collision count.

\section{Explicit analytic and geometric inputs}
We collect the indispensable inherited assertions in the exact form
used by the new proofs. Their source theorems and operator dependencies
are specified here and in Appendix~\ref{app:v83-kernels}. The new
potential and first-exit estimates do not assume a surface spectral
insertion theorem or an unweighted pointwise local bound.

\begin{inputstatement}[Physical flow and source comparison]
\label{input:v83-geometric}
On the regular part $\Omega$ of each exact event, endpoint angles
$q=(\alpha_0,\alpha_m)$ give $F=L$, $g=\nabla_qF=R(-p_0,p_m)$,
$K=D_q^2F$ and $\vartheta=\arg g$. Put
\[
 \lambda_\circ=9\sqrt{19}/200,\qquad \Lambda_\circ=47/6,\qquad
 \omega=|\partial_{p_0}\alpha_m|^{-1},\qquad
 \chi=|\det D_q(F,\vartheta)|.
\]
The bounds $\lambda_\circ I\le K\le\Lambda_\circ I$ and
$\lambda_\circ\le\chi\le\Lambda_\circ$ hold uniformly.
The explicit field
\[
 X_q=\frac{K^{-1}g}{g^{\mathsf T}K^{-1}g}
\]
is stopped at true physical and section-decision boundaries, not at
chart seams. Its flow satisfies $XF=1$, $X\vartheta=0$ and
\begin{equation}\label{eq:v83-radial-input}
 2\lambda_\circ s\le |g(\Phi_sx)|^2-|g(x)|^2
                                      \le2\Lambda_\circ s.
\end{equation}
Its endpoint-curve length is at most $C\sqrt s$. For
$I_m=1+\sum_{i=1}^{m-1}q_0^{d_i}/c_i$ one has
$|X\omega|\le C\omega I_m/|g|$ on regular survivors. At a genuine
grazing exit some internal incidence tends to zero and $\omega$ tends
to zero. Forward interior critical termination cannot occur.

Every starting point with forward lifetime at most $r$ belongs to a
layer of one of the following forms, at an actual collision or flight:
\begin{equation}\label{eq:v83-layer-input}
 c_j\le C_Eq_0^{d_j}\sqrt r,\quad
 \Delta_j\le C_Eq_0^{d'_j}\sqrt r,\quad
 \operatorname{dist}(T_R^jx,\partial Y_R^*)\le C_Eq_0^{d_j}\sqrt r.
\end{equation}
The physical coarea covariance, proved in the appendix in the convention
used here, has Jacobian $J_h=(\omega(\Phi_hx)/\omega(x))\chi(x)/\chi(\Phi_hx)$.
Writing $E_h^\varepsilon$ for all weighted starting exit mass and
\[
 O_h^\varepsilon(t)=\int_{S_h(t)} W^\varepsilon(x)
            [1-\omega(\Phi_hx)/\omega(x)]_+\,d\mu_t(x),
\]
the complete source satisfies, for each fixed $r$ and almost every $t$,
\begin{equation}\label{eq:v83-shift-input}
 b^\varepsilon(t)\le2A_rb^\varepsilon(t)+\bar E_r^\varepsilon(t)
       +\bar O_r^\varepsilon(t)
       +C_\Theta(1+\varepsilon^{-1})\sqrt r\,A_rp(t),
 \qquad A_rf(t)=\frac1{2r}\int_{t-r}^{t+r}f(s)\,ds.
\end{equation}
\end{inputstatement}
These are the endpoint-field, all-contact motion, uniform cubic-action
and separated-volume results in retained modules 174, 177, 180--182.
The guard term in \eqref{eq:v83-shift-input} is the complete product
guard, not a protective subword source. The derivative and continuation
assertions are confined to regular physical trajectories.

\begin{inputstatement}[Complete local source estimates]
\label{input:v83-local}
For every fixed auxiliary band $B_*$, constants $C_{B_*}<\infty$ and
$\rho_{B_*}<1$ give the following bounds uniformly in radius, exact
labels, interval $J$ and mark $0\le j\le m$. Let $S_{R,s}$ be the
union of strips of width $Cs$ about the finitely many coordinate lines
supporting $\partial Y_R^*$. For $s$ small,
\begin{align}
 m^2\nu\{K_m^{\rm c}=k,A_m=n,L_m\in J,
                       T_R^jx\in S_{R,s}\}
 &\le Cs^{1/16}(|J|+B_*^{-1})
   +C_{B_*}(1+B_*|J|)m^2\rho_{B_*}^{m/2},\label{eq:v83-strip-input}\\
 m^2\int_Jb^\delta(t)\,dt
 &\le C\delta^{1/16}(|J|+B_*^{-1})
   +C_{B_*}(1+B_*|J|)m^3\rho_{B_*}^{m/2}.\label{eq:v83-guard-input}
\end{align}
The local extension to a fixed remaining compact range of $\delta$
is included. At auxiliary parameter one $W^1=1$, so $b^1=p$.

Put $z=|g|$, $\alpha=129/128$ and, for $0<\eta<1/2$, define the
actual nonnegative source measures
\begin{align}
 v^{\eta,b}(t)\,dt&=L_\#(I_m z^{-\eta}\one_{\{0<z<b\}}
                                       \one_\Pi\nu_R^*),\notag\\
 v^\eta(t)\,dt&=L_\#(I_m z^{-\eta}\one_\Pi\nu_R^*).
                                                   \label{eq:v83-fractional-input}
\end{align}
They are finite and absolutely continuous at every fixed count, and
for $0<b<b_1$,
\begin{equation}\label{eq:v83-fractional-finite-input}
 m^2\int_Jv^{\eta,b}(t)\,dt
 \le C_\eta b^{\alpha-\eta}(|J|+B_*^{-1})
  +C_{B_*,\eta}(1+B_*|J|)m^2b^{1/2-\eta}\rho_{B_*}^{m/4}.
\end{equation}
Adding the bounded-gradient part gives, for every fixed $L>0$,
\begin{equation}\label{eq:v83-fixed-width-input}
 \limsup_m\sup_{R,n,k}\sup_{|J|=L}m^2\int_Jv^\eta\le C_\eta L,
 \qquad
 \limsup_m\sup_{R,n,k}\sup_{|J|=L}m^2\int_Jv^{\eta,b}
                                      \le C_\eta b^{\alpha-\eta}L.
\end{equation}
Every auxiliary band is fixed before the collision limit and enlarged
only afterwards; leading constants are independent of that band and
of the previously fixed interval length.
\end{inputstatement}
Equation~\eqref{eq:v83-strip-input} is Supplement
Theorem~\ref{supp-thm:v48-marked-strip-local}. The physical source
bound \eqref{eq:v83-guard-input} is Supplement
Theorem~\ref{supp-thm:v49-physical-remainder-local}, in particular its
finite-window estimate. The fractional bounds are the positive
layer-cake conclusions of retained module 187, stated in full above.
The joint spectral and common-space assumptions used there are recorded
in the appendix. No estimate for an untruncated first inverse gradient
or for a singular roof convolution is part of this input.

\begin{inputstatement}[The scalar positive remainder]
\label{input:v83-scalar}
For large fixed $B$, with $\varepsilon(B)=A_0B^{-1/12}$,
\begin{equation}\label{eq:v76-remainder}
 m^2p-\mathcal L=m^2b^{\varepsilon(B)}+e_B,\qquad
 \limsup_m\sup_{R,n,k}\esssup_{\mathcal C(M)}|e_B|
                                      \le C_MB^{-1/192}.
\end{equation}
The source $b^\varepsilon$ is nonnegative and has a compactly supported
$BV$ density at fixed count. No collision-uniform variation bound is
assumed. Supplement Theorem~\ref{supp-thm:v51-positive-error} gives
this remainder with the same complete finite-peak reference.
\end{inputstatement}

\input{core/190_untruncated_optical_potentials}
\input{core/191_complete_first_exit_surface_mass}
\input{core/192_ordered_fractional_tail_and_endpoint}

\section{The scalar objective and downstream normalization}
The organizing assertion remains \eqref{eq:v76-target}. The complete
optical norm and first-exit surface mass are now quantitatively bounded
without an initial critical cutoff. They are not essential-height
bounds, and the two height terms in \eqref{eq:v83-endpoint} remain
distinct proof obligations. In particular a first-exit surface is not
removed by an internal signed occupation difference, and a fixed-width
volume theorem is not used at a collision-dependent microscopic width.

After scalar closure, the positive path-remainder equality in
Supplement Theorem~\ref{supp-thm:v73-scalar-path-equality} supplies
the same-roof path numerator. Ratios are formed only on positive
arithmetic reference classes; the actual-return clock and forward
likelihood retain their stated inherited hypotheses. Nothing in the
new estimates divides by an arithmetic zero class or reverses this
scalar-first order.

\appendix
\input{appendices/v83_common_kernels_and_dependencies}
\input{journal-references}
\end{document}
